\section{Conclusion}
\label{sec:conclusion}

We evaluated a leakage-controlled framework that combines RFECV with a heterogeneous stacking ensemble for T2DM prediction, using repeated cross-validation with corrected inference, an untouched hold-out set, calibration, threshold and decision-curve analysis, and SHAP explanations. RFECV produced compact feature sets built on glucose, BMI, family history and their interactions, without loss of discrimination. The stacking ensemble achieved the highest observed hold-out sensitivity (0.815), F1 (0.710) and MCC (0.532) on this split. However, across 50 cross-validation folds and on the hold-out set, it did not discriminate significantly better than regularised logistic regression (cross-validated ROC-AUC 0.843 versus 0.842). Class weighting produced a systematic overestimation of risk that discrimination metrics do not reveal. On small clinical tabular datasets, increased ensemble complexity therefore does not by itself improve prediction. The main value of such a framework lies in rigorous, leakage-free evaluation, sensitivity-oriented operating points and transparent reporting of calibration and clinical utility.
